\documentclass[letterpaper,11pt]{moderncv}
\moderncvstyle{banking}
\moderncvcolor{black}

\usepackage[letterpaper, margin=0.8in]{geometry}
\usepackage{xcolor}
\usepackage{fontspec}
\setmainfont{Times New Roman}
\usepackage[unicode]{hyperref}

\name{Saeyoung}{Kim}
\address{Cambridge, MA, USA}{}{}
\phone[mobile]{+1 (351) 322 5510}
\email{saeyoung\_kim@uml.edu}
\social[linkedin]{saeyoung-macx-kim}
\social[github]{PrimaryAnomaly}

\recipient{Hiring Team}{Figure\\San Jose, CA}
\date{\today}
\opening{Dear Hiring Team,}
\closing{Sincerely,}

\begin{document}

\makelettertitle

I am applying for the Reliability Test Engineer, Mechatronics position at Figure. I have 9+ years of experience designing and building test stations for electromechanical systems, and I want to build the test infrastructure for your humanoid robots.

At Hanwha Systems, I designed EGSE test stations for satellite subsystem validation, wiring up power supplies, electronic loads, and DAQ instrumentation. I built test interfaces over UART, I2C, SPI, and Ethernet, qualified stations for accuracy and repeatability, and partnered with hardware engineers on accelerated life testing. When sub-assemblies failed, I diagnosed root causes and worked with suppliers to fix the test equipment.

During my PhD, I built scalable test systems for multi-modal sensor suites and wrote embedded DAQ firmware in C/C++. I programmed automated test sequences, designed test interfaces, and did hands-on PCB assembly across many prototype iterations. In my postdoc, I develop test automation frameworks in Python, validate test station accuracy for long-running experiments, and plan lab equipment setups with engineering teams.

Figure is building something that has never existed at scale before, and that means the test infrastructure matters as much as the design. I have the test station design, protocol integration, automation, and debugging experience to help make that happen.

Thank you for your consideration.

\makeletterclosing

\end{document}
